% TOP_PROBLEM: {"id": 30006866, "title": "Sinai's Positive Metric Entropy Conjecture for the Chirikov Standard Map", "category_id": 11, "category": {"id": 11, "name": "dynamical_systems", "display_name": "Dynamical Systems", "description": "Problems about long-term behavior of deterministic systems, Hamiltonian dynamics, and periodic orbits.", "slug": "dynamical-systems", "order_index": 11, "created_at": "2026-07-31T15:26:25.671Z"}, "status": "open"}
% FIRST_POSTED: 2026-09-25
% ENTRY_KIND: statement_only
% !TeX program = lualatex
% Definition and sources only; this file is not an LLM solution attempt.
\documentclass[11pt]{article}
\usepackage[margin=25mm]{geometry}
\usepackage{fontspec}
\setmainfont{Latin Modern Roman}
\usepackage{amsmath,amssymb,mathtools}
\usepackage{unicode-math}
\setmathfont{Latin Modern Math}
\usepackage{xurl,hyperref}
\hypersetup{colorlinks=true,urlcolor=blue}
\setcounter{secnumdepth}{0}
\setlength{\parindent}{0pt}
\setlength{\parskip}{5pt}
\setlength{\emergencystretch}{3em}
\begin{document}
\section*{\texorpdfstring{Sinai's Positive Metric Entropy Conjecture for the Chirikov Standard Map}{Sinai's Positive Metric Entropy Conjecture for the Chirikov Standard Map}}\label{30006866}
\subsection{Definitions and mathematical statement}
On $\mathbb T^2={\mathbb{R}}^2/{\mathbb{Z}}^2$, define the area-preserving standard map
\[
 f_k(x,y)=\bigl(x+y+k\sin(2\pi x),\ y+k\sin(2\pi x)\bigr)\pmod{{\mathbb{Z}}^2}.
\]
Let $m$ be normalized Lebesgue area. Its metric entropy $h_m(f_k)$ measures the asymptotic information rate of measurable partitions under iteration, taking the supremum over finite partitions. The conjectured assertion is
\[
 \operatorname{Leb}\{k\in{\mathbb{R}}:h_m(f_k)>0\}>0.
\]
For this smooth area-preserving map, the entropy formula identifies positivity with a positive-$m$-measure set of points having positive maximal Lyapunov exponent, where that exponent is the long-time exponential growth rate of tangent vectors. Positive topological entropy on a zero-area invariant set is not sufficient. The parameter normalization here is exactly the one in the catalogue.
\par\smallskip\textit{Formulation and concept references:} \href{https://www.math.uci.edu/~asgor/Psim076.pdf}{[S1]}.

\subsection{Short English statement}
For a positive-measure set of parameters, does the standard map exhibit chaotic tangent growth on a positive-area set of trajectories?
\subsection{Sources}
\begin{itemize}
\item[S1]A. Gorodetski and V. Kaloshin, "Conservative Homoclinic Bifurcations and Some Applications," Proc. Steklov Inst. Math. 267 (2009), 76-90. DOI:10.1134/S0081543809040063.
\url{https://www.math.uci.edu/~asgor/Psim076.pdf}.
\textit{Formulation source.}
\end{itemize}
\end{document}
